\section{Background}
\label{sec:background}
% origin
% instructions
% universal basis

\begin{table}
    \centering
    % --- First Sub-table ---
    \begin{subtable}[t]{0.49\linewidth}
        \centering
        \begin{tabular}{cl}
            \toprule
            \texttt{>} & Increment pointer \\
            \texttt{<} & Decrement pointer \\
            \texttt{+} & Increment byte at pointer \\
            \texttt{-} & Decrement byte at pointer \\
            \texttt{.} & Print byte at pointer \\
            \texttt{,} & Write user input to byte at pointer \\
            \texttt{[} & Go past matching \texttt{]} if byte at pointer is 0 \\
            \texttt{]} & Return to matching \texttt{[} if byte at ptr is not 0\\
            \bottomrule
        \end{tabular}
        \caption{\bfck{} commands}
        \label{tab:bfck-commands}
    \end{subtable}
    \begin{subtable}[t]{0.49\linewidth}
        \centering
        \begin{tabular}{cl}
            \toprule
            \texttt{\}} & Increment pointer \\
            \texttt{\{} & Decrement pointer \\
            \texttt{*} & Apply $X$ gate to qubit at pointer \\
            \texttt{\~} & Apply $H$ gate to qubit at pointer \\
            \texttt{;} & Apply $T$ gate to qubit at pointer \\
            \texttt{:} & Measure qubit at pointer and print outcome \\
            \texttt{\#} & Control next gate on the qubit at pointer \\
            \texttt{?} & Apply next gate iff last measurement was 1 \\
            \bottomrule
        \end{tabular}
        \caption{\bfq{} commands}
        \label{tab:bfq-commands}
    \end{subtable}
    
    \caption{This is the main caption for both tables side by side.}
    \label{tab:main}
\end{table}
